\documentclass[10pt]{article}
\usepackage[margin=0.82in]{geometry}
\usepackage[T1]{fontenc}
\usepackage{lmodern}
\usepackage[hidelinks]{hyperref}
\usepackage{enumitem}
\setlength{\parindent}{0pt}
\setlength{\parskip}{0.52em}
\setlist[enumerate]{leftmargin=1.7em,itemsep=0.22em,topsep=0.2em}
\begin{document}

{\Large\bfseries Cover Letter}\par
19 August 2026\par
Editors of \emph{Ledger}\\
University of Pittsburgh

Dear Editors,

I submit the manuscript \textbf{“MOSAIC: State-Transition-Centric Integrity Capsules for Distributed Ledger Execution”} for consideration as a full-length research article in \emph{Ledger}.

The manuscript introduces \emph{Evidence-Complete Transition Closure} (ECTC), a protocol semantics in which a state transition is not terminal merely because a validator accepted it. A Capsule must reach a verifiable \emph{Closed} outcome carrying a ClosureProof and successor StateSeal, a \emph{Conflict} outcome carrying ConflictEvidence, or an \emph{Abandoned} outcome carrying an AbandonProof. The proposed lifecycle binds predecessor state, transition claim, witness obligation, and successor state while preserving incompatible and abandoned claims as independently inspectable artifacts.

The paper makes three contributions. First, it defines and implements the Capsule--WitnessReceipt--first-claim lock--ClosureProof--StateSeal lifecycle and the ECTC outcome classifier. The manuscript now gives an explicit Valid ClosureProof predicate requiring context binding, unique signer/receipt correspondence, individually valid ACCEPT receipts, threshold weight, and recomputation of the canonical proof digest; Closed additionally requires a matching successor StateSeal. Second, it states eight protocol invariants, separates an unweighted BFT quorum model from a weighted extension, and gives a conditional quorum-intersection argument for closure exclusivity. It also defines Evidence Completeness as a state-observation property and Liveness as a separate temporal property, so ECTC is not presented as an eventual-termination guarantee. The corresponding TLA+ model and executable checks are explicitly bounded and are not presented as an unbounded formal proof. Third, it supplies a reproducible artifact with adversarial state-machine tests, disjoint-versus-competing-claims workloads, batching measurements, validator-process scaling, byte decomposition, and crash/recovery evaluation.

The implementation currently collects 102 automated tests. The model checker checks the unweighted intersection property over 84 finite committee cases and 36,703 quorum pairs, plus weighted and ECTC checks. In a 1,000-operation serialized local workload, disjoint and genuine competing-claims cases measured 3.54/3.63 ms and 4.19/5.02 ms p50/p95 latency, respectively; the competing-claims run rejected 1,000 losing claims and generated 1,000 ConflictEvidence artifacts at 635 bytes per rejected claim on average. Batch sizes of 10 and 100 measured 2,530.78 and 2,473.96 protocol bytes per successful operation. A seven-process local rehearsal with two Byzantine test identities, a kill/restart, and 24 partial-frame probes completed all 120 operations with liveness 1.0, zero unexpected operational errors, p50 36.80 ms, and p95 233.56 ms.

The manuscript deliberately labels these results \texttt{LOCAL\_EMULATION}. They are not Internet-WAN measurements, permissionless Sybil-resistance evidence, production-mainnet evidence, or a claim of superiority over HotStuff, Narwhal/Tusk, FastPay, Sui Lutris, Mysticeti, Block-STM, or BFT-object systems. The related-work section treats object-centric execution, optimistic concurrency, selective ordering, certificates, and BFT CRDTs as prior art. The scientific claim is narrower and falsifiable: ECTC is a semantic interface that makes terminal positive and negative transition outcomes explicit and auditable, while preserving Pending as non-terminal; it does not claim priority over the individual mechanisms it composes.

The public research artifact is available at:

\url{https://github.com/adnanomar77/mosaic-protocol}

It contains the source, tests, formal model, benchmark drivers, adversarial matrix, Mermaid sources, rendered figures, CSV data, JSON results, and checksums. Code is released under Apache-2.0; manuscript and documentation are released under CC BY 4.0 where applicable. The repository excludes private keys, credentials, virtual environments, and runtime WAL data.

An AI-assisted software and language workflow was used for code inspection, drafting support, organization, and language editing. The author ran the experiments, reviewed the implementation, checked the sources and reported values, and remains solely responsible for the manuscript, claims, limitations, and final submission. This use is disclosed in the manuscript in accordance with \emph{Ledger}'s AI policy.

The manuscript is original, is not under consideration elsewhere, and has not been published in a peer-reviewed venue. No external funding was received. The author declares no known conflicts of interest relevant to this submission.

\textbf{Suggested reviewers.} The following scholars are suggested because their public institutional profiles document expertise directly relevant to distributed consensus, blockchain protocols, cryptographic security, and reproducible distributed systems. The author has no known personal, institutional, supervisory, employment, or recent co-authorship relationship with any of them. The editor should exclude any candidate for whom an undisclosed conflict or availability issue is identified; none has been contacted or has agreed to review.

\begin{enumerate}
\item \textbf{Christian Cachin} --- Professor of Computer Science, University of Bern. Expertise: blockchain and consensus protocols, distributed computing, cryptographic protocols, and cloud-computing security. Public institutional email: \texttt{christian.cachin@unibe.ch}. Profile: \url{https://crypto.unibe.ch/cc/}
\item \textbf{Rachid Guerraoui} --- Full Professor, Distributed Computing Laboratory, EPFL. Expertise: distributed algorithms, secure distributed storage, transactional shared memory, distributed programming languages, and robust distributed computing. Public institutional email: \texttt{rachid.guerraoui@epfl.ch}. Profile: \url{https://people.epfl.ch/rachid.guerraoui?lang=en}
\item \textbf{Elaine Shi} --- Professor of Computer Science and Electrical and Computer Engineering, Carnegie Mellon University. Expertise: blockchain, cryptography, distributed systems security, formal methods, protocol security, and verification. Public institutional email: \texttt{rshi@andrew.cmu.edu}. Profile: \url{https://www.cylab.cmu.edu/directory/bios/shi-elaine.html}
\end{enumerate}

Thank you for considering this submission. I would welcome a rigorous review of both the protocol hypothesis and the limits of the current evidence.

Sincerely,\\[0.15em]
\textbf{Adnan Omar Awad Allemon}\\
Independent Researcher\\
Corresponding author: \href{mailto:adnanomar774@gmail.com}{adnanomar774@gmail.com}

\end{document}
